% Part: first-order-logic
% Chapter: syntax-and-semantics
% Section: first-order-languages

\documentclass[../../../include/open-logic-section]{subfiles}

\begin{document}

\olfileid[pt]{fol}{syn}{fol}

\olsection{Linguagens de Primeira Ordem}


As expressões da lógica de primeira ordem são construídas a partir de um vocabulário básico
contendo \emph{!!{variable}s}, \emph{!!{constant}s},
\emph{!!{predicate}s} e às vezes \emph{!!{function}s}.  A partir deles,
juntamente com conectivos lógicos, quantificadores e símbolos de pontuação
como parênteses e vírgulas, formam-se \emph{termos} e
\emph{!!{formula}s}.

\begin{explain}
Informalmente, !!{predicate}s são nomes para propriedades e relações,
!!{constant}s são nomes para objetos individuais, e !!{function}s são
nomes para mapeamentos.  Estes, exceto a !!{identity}~$\eq$, são os
\emph{símbolos não lógicos} e juntos constituem uma linguagem.  Qualquer
linguagem de primeira ordem~$\Lang L$ é determinada por seus símbolos não
lógicos.  No caso mais geral, $\Lang L$ contém infinitos
símbolos de cada tipo.
\end{explain}

No caso geral, fazemos uso dos seguintes símbolos na
lógica de primeira ordem:

\begin{enumerate}
\item Símbolos lógicos
\begin{enumerate}
\item Conectivos lógicos:
  \startycommalist
  \iftag{prvNot}{\ycomma $\lnot$ (negação)}{}%
  \iftag{prvAnd}{\ycomma $\land$ (conjunção)}{}%
  \iftag{prvOr}{\ycomma $\lor$ (disjunção)}{}%
  \iftag{prvIf}{\ycomma $\lif$ (!!{conditional})}{}%
  \iftag{prvIff}{\ycomma $\liff$ (!!{biconditional})}{}%
  \iftag{prvAll}{\ycomma $\lforall$ (quantificador universal)}{}%
  \iftag{prvEx}{\ycomma $\lexists$ (quantificador existencial)}{}.
\tagitem{prvFalse}{A constante proposicional para !!{falsity}~$\lfalse$.}{}
\tagitem{prvTrue}{A constante proposicional para !!{truth}~$\ltrue$.}{}
\item A !!{identity} de dois lugares~$\eq$.
\item Um conjunto !!{denumerable} de !!{variable}s: $\Obj v_0$, $\Obj v_1$, $\Obj
  v_2$, \dots
\end{enumerate}
\item Símbolos não lógicos, que constituem a \emph{linguagem
  padrão} da lógica de primeira ordem
\begin{enumerate}
\item Um conjunto !!{denumerable} de !!{predicate}s de $n$ lugares para cada $n>0$: $\Obj
  A^n_0$, $\Obj A^n_1$, $\Obj A^n_2$, \dots
\item Um conjunto !!{denumerable} de !!{constant}s: $\Obj c_0$, $\Obj c_1$, $\Obj
  c_2$, \dots.
\item Um conjunto !!{denumerable} de !!{function}s de $n$ lugares para cada $n>0$:
  $\Obj f^n_0$, $\Obj f^n_1$, $\Obj f^n_2$, \dots
\end{enumerate}
\item Sinais de pontuação: (, ) e a vírgula.
\end{enumerate}

A maioria de nossas definições e resultados será formulada para a linguagem padrão
completa da lógica de primeira ordem.  No entanto, dependendo da
aplicação, também podemos restringir a linguagem a apenas alguns
!!{predicate}s, !!{constant}s e !!{function}s.

\begin{ex}
A linguagem~$\Lang L_A$ da aritmética contém um único !!{predicate} de dois lugares~$<$,
um único !!{constant}~$\Obj 0$, um !!{function} de um lugar~$\prime$,
e dois !!{function}s de dois lugares~$+$ e~$\times$.
\end{ex}

\begin{ex}
A linguagem da teoria dos conjuntos~$\Lang L_Z$ contém apenas o único
!!{predicate} de dois lugares~$\in$.
\end{ex}

\begin{ex}
A linguagem das ordens~$\Lang L_\le$ contém apenas o !!{predicate}
de dois lugares~$\le$.
\end{ex}

Novamente, estas são convenções: oficialmente, estes são apenas apelidos,
por exemplo, $<$, $\in$ e $\le$ são apelidos para $\Obj A^2_0$, $\Obj 0$ para
$\Obj c_0$, $\prime$ para $\Obj f^1_0$, $+$ para $\Obj f^2_0$, $\times$ para
$\Obj f^2_1$.

\iftag{defNot,defOr,defAnd,defIf,defIff,defTrue,defFalse,defEx,defAll}{%
Além dos conectivos primitivos e
\iftag{notprvEx,notprvAll}{quantificador}{quantificadores} introduzidos
acima, também usamos os seguintes símbolos \emph{definidos}:
\startycommalist
  \iftag{defNot}{\ycomma $\lnot$ (negação)}{}%
  \iftag{defAnd}{\ycomma $\land$ (conjunção)}{}%
  \iftag{defOr}{\ycomma $\lor$ (disjunção)}{}%
  \iftag{defIf}{\ycomma $\lif$ (!!{conditional})}{}%
  \iftag{defIff}{\ycomma $\liff$ (!!{biconditional})}{}%
  \iftag{defAll}{\ycomma $\lforall$ (quantificador universal)}{}%
  \iftag{defEx}{\ycomma $\lexists$ (quantificador existencial)}{}%
  \iftag{defFalse}{\ycomma !!{falsity}~$\lfalse$}{}%
  \iftag{defTrue}{\ycomma !!{truth}~$\ltrue$}}{}.

\begin{tagblock}{defNot,defOr,defAnd,defIf,defIff,defTrue,defFalse,defEx,defAll}
\begin{explain}
Um símbolo definido não faz parte oficialmente da linguagem, mas é
introduzido como uma abreviação informal: ele nos permite abreviar
fórmulas que, se usássemos apenas símbolos primitivos, ficariam bastante
longas.  Isso é obviamente uma vantagem.  A maior vantagem, no entanto,
é que as provas ficam mais curtas.  Se um símbolo é primitivo, ele tem que ser
tratado separadamente nas provas. Quanto mais símbolos primitivos, portanto,
mais longas nossas provas.
\end{explain}
\end{tagblock}

% Alternate symbols

\begin{intro}
Você pode estar familiarizado com terminologia e símbolos diferentes dos
que usamos acima. Textos de lógica (e professores) comumente usam 
$\sim$, $\neg$, ou~!\ para ``negação'', $\wedge$, $\cdot$, ou $\&$
para ``conjunção''.  Símbolos comumente usados para o ``condicional'' ou
``implicação'' são $\rightarrow$, $\Rightarrow$, e $\supset$.
\iftag{prvIff,defIff}{Símbolos para ``bicondicional'', ``bi-implicação'',
  ou ``equivalência (material)'' são $\leftrightarrow$,
  $\Leftrightarrow$, e $\equiv$.}{}
\iftag{prvFalse,defFalse}{O símbolo $\lfalse$ é chamado de várias formas:
  ``falsidade'', ``falsum'', ``absurdo'', ou ``bottom''.}{}
\iftag{prvTrue,defTrue}{O símbolo $\ltrue$ é chamado de várias formas:
  ``verdade'', ``verum'', ou ``top''.}{}

É convencional usar letras minúsculas (por exemplo, $a$, $b$, $c$) do
início do alfabeto latino para !!{constant}s (às vezes chamados
de nomes), e letras minúsculas do final (por exemplo, $x$, $y$, $z$) para
!!{variable}s. Quantificadores combinam-se com !!{variable}s, por exemplo, $x$;
variações notacionais incluem $\forall x$, $(\forall x)$, $(x)$, $\Pi x$,
$\bigwedge_x$ para o quantificador universal e $\exists x$, $(\exists
x)$, $(Ex)$, $\Sigma x$, $\bigvee_x$ para o quantificador existencial.
\end{intro}

\begin{explain}
Poderíamos tratar todos os operadores proposicionais e ambos os quantificadores como
símbolos primitivos da linguagem.  Poderíamos, em vez disso, escolher um estoque
menor de símbolos primitivos e tratar os outros !!{operator}s como
definidos. Conjuntos ``funcionalmente completos para a verdade'' de operadores booleanos
incluem $\{ \lnot, \lor \}$, $\{ \lnot, \land \}$, e $\{ \lnot,
\lif\}$---estes podem ser combinados com qualquer um dos quantificadores para uma
linguagem de primeira ordem expressivamente completa.

Você pode estar familiarizado com dois outros !!{operator}s: a barra de Sheffer~$|$ (em homenagem a Henry Sheffer), e a
seta de Peirce~$\downarrow$, também conhecida como adaga de Quine.  Quando recebem suas
leituras usuais de ``nand'' e ``nor'' (respectivamente), esses operadores
são funcionalmente completos para a verdade por si sós.
\end{explain}

\end{document}
